\chapter*{Abstract}
\addcontentsline{toc}{chapter}{Abstract}

Modern software depends heavily on open-source components, yet downstream
projects remain exposed when an upstream repository quietly stops being
maintained. This thesis asks to what extent publicly observable repository
health and maintenance indicators can predict whether an open-source dependency
becomes inactive within a twelve-month horizon, and how the resulting risk score
can support risk-based dependency triage. Following a Design Science Research
approach, the work develops and evaluates OSS Risk Radar, an artifact that
constructs a leakage-controlled historical dataset from GH Archive event data
and public package metadata, engineers observation-time features describing
activity, responsiveness, and how narrowly a project's work rests on individual
contributors, and trains calibrated inactivity-risk models.

On a repository foundation dataset of 5{,}000 projects and 50{,}000 quarterly
observation snapshots, a gradient-boosted model reaches a held-out ROC-AUC of
\(0.958\) --- the probability that an inactive repository is ranked above an
active one --- with a Brier score of \(0.078\), and its skill largely persists
under a repository-disjoint split (\(0.950\)). This aggregate figure is, however, a poor
measure of what the model contributes: a trivial rule that ranks repositories by
days since their last commit reaches \(0.957\) on the same rows, against
\(0.974\) for the model. The learned model earns its complexity on the
repositories that are still active at the observation time --- the subgroup for
which an early warning is most decision-relevant --- where it attains \(0.905\)
against \(0.844\) for the best trivial rule. The thesis therefore reports both a
null result on the aggregate comparison and a positive result on early decline.
Excluding bot accounts from the commit and contributor signals is shown to be
worth its cost by retraining the model on the unfiltered signals and comparing
both against the same clean labels, and security-practice signals are kept as
parallel triage context rather than as predictors of inactivity. A
practical demonstrator applies the deployed model per repository to the
dependency repositories of a real project, producing a ranked triage view in
which each repository's calibrated score is accompanied by an operational
evidence-support measure. The demonstrator shows how the score can be integrated
into a dependency-triage workflow rather than establishing that it improves
triage decisions. The score is positioned as a decision-support signal for
prioritizing dependency review, not as a verdict on abandonment.
